\documentclass[11pt]{article}

\usepackage[margin=1in]{geometry}
\usepackage{amsmath,amssymb}
\usepackage{booktabs}
\usepackage{graphicx}
\usepackage{siunitx}
\usepackage[hidelinks]{hyperref}

\title{A Closed-Form Conservation-Correction Layer for\\
Error-Bounded Lossy Compression of Scientific Fields}
\author{Alika M. Parks\\
\small Independent Researcher \\
\small \texttt{alikamp@gmail.com}}
\date{September 2026}

\begin{document}
\maketitle

\begin{abstract}
Error-bounded lossy compressors such as ZFP and SZ guarantee a pointwise error
bound but do not preserve integrated quantities---total energy, mass,
momentum---so downstream analyses that check a conservation law observe a bias
that grows with the compression ratio. We describe a closed-form correction
layer that restores exact conservation of a chosen global quadratic invariant
on top of any lossy backend. After the backend reconstructs a field, a single
multiplicative factor is computed in closed form so that the reconstruction's
total energy equals the original's; the factor is stored (\SI{4}{\byte} per
field) and applied on decode. The correction leaves the backend's compression
ratio and pointwise error bound effectively unchanged. On a synthetic
three-dimensional turbulence proxy, applied to ZFP across compression ratios
from $14\times$ to $60\times$, the layer reduces the relative energy-conservation
error by four to six orders of magnitude (from $4.8\times10^{-4}$ to
$7\times10^{-9}$ at $60\times$), and the reduction increases with the
compression ratio. We give the derivation, the extension to multiple
simultaneous invariants, and the method's limitations. Reference code is
released under the Apache-2.0 license.
\end{abstract}

\section{Introduction}
The output volume of numerical simulations and instruments has made lossy
floating-point compression a standard element of scientific data pipelines.
Error-bounded compressors such as ZFP~\cite{zfp} and SZ~\cite{sz} provide a
guaranteed pointwise (per-element) error bound, which controls the local
fidelity of a reconstructed field. They do not, however, provide any guarantee
on \emph{integrated} quantities computed from that field. Many downstream
scientific diagnostics depend on such integrals: total kinetic energy, mass,
and momentum are conservation checks, and their drift under compression is a
systematic bias rather than bounded noise.

This work addresses that gap with a minimal construction. Given any lossy
backend, we compute in closed form the single scale factor that makes the
reconstruction's total energy match the original, store that factor, and apply
it on decode. The contribution is deliberately narrow:
\begin{itemize}
  \item it is not a new codec and does not improve the compression ratio;
  \item it is a backend-agnostic post-processing layer, demonstrated here on
        ZFP;
  \item it converts a chosen global quadratic invariant from
        ``approximately preserved'' to ``exact to floating-point precision,''
        at a storage cost of \SI{4}{\byte} per field and $O(N)$ arithmetic.
\end{itemize}

\section{Related work}
Error-bounded lossy compression for scientific data is an active area; ZFP
\cite{zfp}, SZ~\cite{sz}, and TTHRESH~\cite{tthresh} are representative and
differ in transform and error-control strategy. Prior work has also considered
preservation of derived quantities and quantities of interest during
compression. The present layer is positioned as the simplest form of that idea:
a closed-form, backend-agnostic correction of a single global quadratic
invariant applied as post-processing.
% TODO (do not submit without literature verification): cite the
% quantity-of-interest / constraint-preserving lossy-compression literature
% and state precisely what those methods already provide, so the contribution
% here is scoped against them rather than presented as first-of-kind.

\section{Method}
Let $f\in\mathbb{R}^N$ be a scalar field and let $\hat f = C(f;\tau)$ be its
reconstruction under a lossy backend $C$ at tolerance $\tau$. Define the energy
functional
\begin{equation}
E(x) \;=\; \sum_{i=1}^{N} x_i^2,
\end{equation}
which is proportional to kinetic energy for a velocity component. We seek a
scalar $\alpha$ such that $E(\alpha\hat f) = E(f)$. Since
$E(\alpha\hat f)=\alpha^2 E(\hat f)$, the solution is closed form:
\begin{equation}
\alpha \;=\; \sqrt{\frac{E(f)}{E(\hat f)}}\,.
\label{eq:alpha}
\end{equation}
The encoder stores $\alpha$ alongside the backend stream; the decoder returns
$\alpha\hat f$. By construction $E(\alpha\hat f)=E(f)$ to floating-point
precision, independent of $\tau$ and of the backend.

\paragraph{Effect on the pointwise bound.}
Write $\alpha = 1+\varepsilon$. From \eqref{eq:alpha},
$\varepsilon = \tfrac{1}{2}\,\big(E(f)-E(\hat f)\big)/E(\hat f) + O(\varepsilon^2)$,
so $|\varepsilon|$ is of the order of the backend's relative energy error, which
is small. The corrected pointwise error satisfies
\begin{equation}
\lVert f-\alpha\hat f\rVert_\infty
\;\le\; \lVert f-\hat f\rVert_\infty + |\varepsilon|\,\lVert \hat f\rVert_\infty,
\end{equation}
i.e.\ the backend bound is preserved up to a bounded, small additive term.

\paragraph{Vector fields.}
For a velocity field with components $(u,v,w)$, a per-component factor
$\alpha_c$ conserves each component energy $E(f_c)$ and therefore the total
kinetic energy $\tfrac{1}{2}\sum_c E(f_c)$. The per-field overhead is
\SI{4}{\byte} per component.

\paragraph{Multiple invariants.}
Two invariants can be enforced jointly. To preserve both mass $M(x)=\sum_i x_i$
and energy $E(x)$, apply an affine map $x' = a\,\hat f + b\mathbf{1}$ and solve
the two constraints $M(x')=M(f)$, $E(x')=E(f)$ for $(a,b)$; this reduces to a
scalar quadratic. Energy together with vector momentum leads to a small
constrained least-squares problem. These extensions are stated here as
directions and are not evaluated below.

\section{Experimental setup}
The test field is a synthetic turbulence proxy on a $64^3$ grid: a Taylor--Green
base flow superposed with multi-octave sinusoidal noise and a small number of
localized Gaussian bursts, stored as \texttt{float32}, evaluated over three
random seeds. The backend is ZFP (via the \texttt{zfpy} binding) in
fixed-accuracy mode. For each configuration we report the compression ratio
against the raw \texttt{float32} array, the maximum absolute pointwise error,
and the relative energy drift $|E(\hat f)-E(f)|/E(f)$, with and without the
correction layer. The field is a proxy and not a solution of the Navier--Stokes
equations; results on measured simulation data are left to future work.

\section{Results}
Table~\ref{tab:main} reports the mean over three seeds. Applying the correction
adds \SI{4}{\byte} per field, which does not change the reported ratios at this
grid size. The maximum pointwise error is unchanged to the reported precision.
The relative energy drift falls from the backend's $10^{-4}$--$10^{-6}$ range to
the $10^{-9}$--$10^{-10}$ range. The reduction is larger at higher compression
ratio because the corrected field's energy is pinned to arithmetic precision
regardless of the backend's drift, whereas the backend's own drift increases as
$\tau$ grows.

\begin{table}[h]
\centering
\begin{tabular}{rrrrr}
\toprule
ZFP tol. & ratio & max abs.\ error & ZFP energy drift & + layer energy drift \\
\midrule
0.02 & $14.1\times$ & $2.5\times10^{-3}$ & $5.5\times10^{-6}$ & $3.4\times10^{-10}$ \\
0.1  & $22.2\times$ & $9.3\times10^{-3}$ & $2.5\times10^{-5}$ & $2.9\times10^{-10}$ \\
0.3  & $35.9\times$ & $3.3\times10^{-2}$ & $7.1\times10^{-5}$ & $1.3\times10^{-10}$ \\
1.0  & $59.7\times$ & $1.1\times10^{-1}$ & $4.8\times10^{-4}$ & $7.1\times10^{-9}$  \\
\bottomrule
\end{tabular}
\caption{Compression ratio, maximum absolute pointwise error, and relative
energy drift for ZFP with and without the correction layer, on a $64^3$
synthetic turbulence proxy (mean over three seeds).}
\label{tab:main}
\end{table}

\section{Applications}
Two settings motivate the layer. The first is archival and cloud storage of
simulation output. Aggressive lossy compression reduces storage and egress
cost, but conservation-sensitive workflows often avoid it because the resulting
bias in integrated quantities is uncontrolled. The layer removes that
objection for one class of invariant at a fixed, negligible cost, so a
conservation check run on the decompressed field returns the original value.
The second is in-loop use inside memory-bandwidth-bound solvers. Because the
correction touches only decompressed arrays, it can be applied within a solver
step rather than only to output files. Lattice-Boltzmann methods are a
representative case: their throughput is limited by the traffic of streaming
distribution functions, so compressing that traffic while holding the physical
invariant fixed is a route to larger single-node domains. This direction is
noted here and not benchmarked.

\section{Limitations}
The evaluation uses a synthetic proxy on a single grid size; measured
simulation data and a standard reduction benchmark (e.g.\ SDRBench~\cite{sdrbench})
are not yet included. The layer enforces one quadratic invariant (energy) and
does not improve non-enforced invariants: a global energy rescale perturbs
linear sums, so momentum drift is unchanged or slightly increased. The
correction is a single global scalar; it fixes the aggregate, not the spatial
distribution of error. Results are shown for one backend (ZFP); backend
independence follows from the construction but is demonstrated once. The
multi-invariant extension of Section~3 is derived but not evaluated.

\section{Conclusion}
A closed-form scalar correction makes an error-bounded lossy compressor conserve
a chosen global quadratic invariant exactly, at \SI{4}{\byte} per field and no
change to the compression ratio or the pointwise error bound. On a synthetic
turbulence proxy the relative energy drift of ZFP is reduced by four to six
orders of magnitude, with the reduction growing as compression becomes more
aggressive. The immediate next steps are evaluation on measured data and a
standard benchmark, comparison against quantity-of-interest--preserving methods,
and evaluation of the multi-invariant extension.

\section*{Code and data availability}
Reference code and the benchmark used here are available at
\url{https://github.com/alikamp/conserved-field-compression} under the
Apache-2.0 license; the reported numbers reproduce from the included notebook.

\begin{thebibliography}{9}
\bibitem{zfp}
P.~Lindstrom.
\newblock Fixed-rate compressed floating-point arrays.
\newblock \emph{IEEE Transactions on Visualization and Computer Graphics},
20(12):2674--2683, 2014.

\bibitem{sz}
S.~Di and F.~Cappello.
\newblock Fast error-bounded lossy HPC data compression with SZ.
\newblock In \emph{IEEE International Parallel and Distributed Processing
Symposium (IPDPS)}, pages 730--739, 2016.

\bibitem{tthresh}
R.~Ballester-Ripoll, P.~Lindstrom, and R.~Pajarola.
\newblock TTHRESH: Tensor compression for multidimensional visual data.
\newblock \emph{IEEE Transactions on Visualization and Computer Graphics},
26(9):2891--2903, 2020.

\bibitem{sdrbench}
K.~Zhao, S.~Di, X.~Liang, S.~Li, D.~Tao, Z.~Chen, and F.~Cappello.
\newblock SDRBench: Scientific data reduction benchmark for lossy compressors.
\newblock In \emph{IEEE International Conference on Big Data}, 2020.
% TODO: verify all bibliographic details before submission, and add
% quantity-of-interest / constraint-preserving compression references.
\end{thebibliography}

\end{document}
